% from ega4-20-fr.tex
\oldpage[IV-4]{226}
Au \textsection20, nous introduisons plusieurs variantes de la notion d'application rationnelle, déjà étudiée dans \sref[I]{section:I.7} d'un point de vue encore assez proche du point de vue classique, et pour cette raison assez mal adaptée au cas des préschémas non nécessairement réduits.
Les notions et résultats du \textsection20 sont utilisés pour développer au \textsection21 (n\textsuperscript{os} \sref{IV.21.1} à \sref{IV.21.7}) la notion générale de diviseur et ses propriétés les plus élémentaires.
Cette notion est surtout commode lorsque les anneaux locaux des préschémas considérés sont noethériens et intégralement clos, et surtout lorsqu'ils sont en outre \emph{factoriels} (\sref{IV.21.6} et \sref{IV.21.7}), en raison de son identification dans ce dernier cas avec la notion de \emph{cycle $1$-codimensionnel} (combinaison linéaire de sous-préschémas irréductibles de codimension~$1$).
Dans \sref{IV.21.9}, on détermine les diviseurs sur un préschéma noethérien de dimension~$1$ mais non nécessairement normal, ce qui est utile pour diverses applications.
Les n\textsuperscript{os} \sref{IV.21.11} et \sref{IV.21.12} donnent deux théorèmes importants, dus respectivement à Auslander--Buchsbaum et Van der Waerden, et se rattachant à la notion d'anneau factoriel (les n\textsuperscript{os} \sref{IV.21.9}, \sref{IV.21.11} et \sref{IV.21.12} sont indépendants les uns des autres).
Dans les n\textsuperscript{os} \sref{IV.21.13} et \sref{IV.21.14}, également indépendants des trois précédents, nous étudions une variante utile de la notion d'anneau local factoriel, celle d'anneau local \emph{parafactoriel}, qui s'introduit notamment~\cite{IV-41} dans le développement de théorèmes de comparaison du groupe de Picard d'un préschéma projectif $X$ sur un corps $k$ et d'une « section hyperplane ».
On verra dans \sref{IV.21.14.1} (théorème de Ramanujam--Samuel) que les anneaux locaux parafactoriels sont bien plus nombreux que l'on ne pouvait s'y attendre \emph{a priori}.

Dans \sref{IV.20.5}, \sref{IV.20.6} et \sref{IV.21.15}, nous reprenons les notions précédentes mais d'un point de vue « relatif » à un préschéma de base fixé.
Pour l'instant ces notions ne sont utilisées qu'assez rarement ; en particulier la notion de diviseur relatif n'est guère employée que lorsqu'il s'agit de diviseurs positifs, et dans ce cas elle s'explicite avantageusement sans le secours de la notion de fonction méromorphe relative, à l'aide de la notion d'immersion transversalement régulière de codimension~$1$.
On aura donc avantage à omettre ces sections en première lecture.

\subsection{Fonctions méromorphes}
\EGAReaderTarget{ega4.c04.subsec.20-1}{IV.20.1}

\begin{env}[20.1.1]
\EGAReaderTarget{ega4.c04.statement.20-1-1}{IV.20.1.1}
Soit $(X,\sh{O}_X)$ un espace annelé, et soit $\sh{S}$ un sous-faisceau \emph{d'ensembles} de $\sh{O}_X$.
Pour tout ouvert $U$ de $X$, considérons l'\emph{anneau de fractions} $\Gamma(U,\sh{O}_X)[\Gamma(U,\sh{S})^{-1}]$ (Bourbaki, \emph{Alg. comm.}, chap.~II, \textsection2, n\textsuperscript{o}1).
Il est immédiat que l'application $U\mapsto\Gamma(U,\sh{O}_X)[\Gamma(U,\sh{S})^{-1}]$ est un \emph{préfaisceau d'anneaux} (\sref[0]{0.1.5.1} et \sref[0]{0.1.5.7}).
On note $\sh{O}_X[\sh{S}^{-1}]$ le \emph{faisceau d'anneaux} associé à ce préfaisceau et on dit que c'est le \emph{faisceau d'anneaux de fractions de $\sh{O}_X$, à dénominateurs dans $\sh{S}$} ; c'est un $\sh{O}_X$-module \emph{plat}.
Il est immédiat que pour tout $x\in X$, on a un isomorphisme canonique
\[
\label{IV.20.1.1.1}
  (\sh{O}_X[\sh{S}^{-1}])_x\isoto\sh{O}_{X,x}[\sh{S}_x^{-1}],
  \tag{20.1.1.1}
\]
car le raisonnement de \sref[0]{0.1.4.5} se généralise aussitôt au cas où l'on a un système inductif $(A_\alpha,\vphi_{\beta\alpha})$ d'anneaux, et pour chaque indice $\alpha$ une partie $S_\alpha$ de $A_\alpha$ telle que
